\scanpage{43}
% archive foliation: 42
Si $u$ nilpotent, il f et s. que $\exp u$ soit un autom d'algèbre.

Op. sur $\operatorname{Hom}_K(M,N)$, quand $\overline G_a$ (resp. $G_a$) opère sur $M$ par $\bar e_u$ (resp. $e_u$), sur $N$ par $\bar e_v$ (resp. $e_v$). Alors il opère sur $\operatorname{Hom}_K(M,N)$ par $\bar e_\delta$ $(e_\delta)$ où
\[
\delta(f)=vf-fu
\]
Mais en \unsure{tte|tt} \uline{rigueur} $\bar e_\delta$ et $e_\delta$ ne sont définis ici que si $\operatorname{Hom}_K(M,N)$ est \inserted{\unsure{is. à}} un foncteur de la forme $W(M)$, ce qui n'est guère assuré que si $M$ proj. de t.f.

Mais on peut généraliser théorie de l'exponentielle aux \deleted{\unsure{tous les}} endomorphismes de tt $K$-foncteur $\mathcal M$ en $\unsure{\mathcal O|\mathbb{Q}}$-modules
\diagram{\begin{tikzcd}[row sep=large,column sep=large]
\mathcal M(K')\arrow[d,"\text{est }(K'\to K'')\text{-semi-linéaire}"]&\text{est un }K'\text{-module}\\
\mathcal M(K'')&
\end{tikzcd}}
Alors $\operatorname{End}_K(\mathcal M)$ est une $K$-algèbre.

a) Si $u\in N(\operatorname{End}_K(\mathcal M))$, $\exp u\in\operatorname{End}_K(\mathcal M)$

et aussi $\exp(\lambda u)$ $\forall\lambda\in G_{a\unsure{K}}(K')\simeq K'$ défini $\in\operatorname{End}_{K'}(\mathcal M_{K'})$

On trouve sur $\mathcal M$ une structure de $G_a$-Module \unsure{lin.|les}
